\begin{abstract}
Open-source tools for map-based LiDAR localization of mobile robots are still largely organized by sensor suite: a 2D scanner calls for an AMCL-style particle filter, a 3D LiDAR for an NDT-based tracker, and an inertial- or GNSS-aided platform for a Kalman fusion node, each with its own topics, parameters, initialization, and recovery behavior. Changing the sensor therefore means re-integrating the localization layer. We describe \emph{\sname}, an open-source ROS~2 localization stack that places three established estimator families behind one node contract: \code{laser2d}, a likelihood-field Monte Carlo localizer; \code{ndt3d}, a planar NDT-MCL backend \citep{saarinen2013ndtmcl} that reuses the same particle filter; and \code{fusion3d}, a 15-state error-state Kalman filter fusing LiDAR, IMU, and RTK-GNSS. Every backend publishes the same \code{map}$\to$\code{odom} transform and pose with covariance and accepts the same \code{/initialpose} seeding, and all filter mathematics live in Eigen-only C++17 cores tested without a ROS runtime. In addition, \code{laser2d} relocalizes a kidnapped robot by deterministic branch-and-bound scan matching over a map-sized window, keeps the top-$K$ candidate poses, commits one only after several odometry-chained scans concentrate a normalized evidence score on it, and otherwise reports the place as ambiguous. In seeded closed-loop experiments on \emph{synthetic} data that drive the actual cores, on 40\,m and 80\,m maps with repeated structure the single-scan rule accepts a wrong pose on 6 and 15 of 120 queries, which multi-scan verification reduces to 0 and 1 of 120 (pooled 95\% Wilson interval 0.1--2.3\%) while committing as many correct poses (242 against 241 of 300). A 100-run NEES test shows that the ESKF covariance is overconfident (full-state NEES above the 95\% interval at 89 of 120 epochs), so \code{fusion3d} is experimental. The contribution is this integration, released as software. No external baseline and no real or public data are used, and no field-accuracy claim is made. Code: \url{https://github.com/kjungmo/prism_loc}.
\end{abstract}
